% Propietario reproducido por unidad demostrativa; véase PROPIETARIOS_CONTINUO_UNIVERSO.json.
% Origen: XI ES CUMPLIMIENTO_EDITORIAL_20260928, sections/ch_cierre_cardinal.tex:26--403.
% Cuerpos y pruebas conservados; sólo namespace, rutas y remisiones contextuales declaradas.
\subsection{Prolongement d’états et extension sémantique transfinie}

Soit

\[
 \Sigma_{\rm HMT}
 =\bigl(\dot R_{\rm APP},\dot R_{\rm TRIT},\dot R_{\rm TPK},
 \dot R_{U},\dot R_{\Gamma_9},\dot R_{\rm Ext},\dot R_\rho,
 \dot R_{\pi},\dot R_{\varphi},\dot R_e,\dot R_{\alpha},
 \dot R_{\rm ar},\dot R_{\rm exc}\bigr)
\tag{33}\label{hmtfg:base:eq:ch-signatura}
\]

la signature \textbf{relationnelle} dénombrable dont le code effectif est \(h\). Chaque \(\dot R_S\) est le graphe de l'opérateur \(S\) sur les codes où il est défini ; on ne présuppose pas qu'une fonction externe soit totale dans chaque niveau. Les graphes \(\dot R_{\pi},\dot R_{\varphi},\dot R_e,\dot R_{\alpha}\) sont les lecteurs de sortie construits par la génération coinductive et la fermeture dodécaphasique ; ils ne reçoivent pas comme arguments les valeurs conventionnelles des constantes. Leur présence dans \(\Sigma_{\rm HMT}\) assure que la clôture conserve l'état complet : elle intègre simultanément les routes exprimant une valeur et l'incidence qui se prolonge vers la réalisation exceptionnelle.

À chaque horizon, le code de l'état enregistre les coordonnées APP--TRIT--TPK, la route, le feuillet, la retenue, la frontière et la mémoire ; à ses étapes de représentation, il enregistre également \(P_N,\Phi_N,E_N,A_N\), où \(A_N\) est le préfixe de \(\alpha_{\rm HMT}\) produit par le même état. Ces coordonnées ne sélectionnent pas la branche de l'extérieur : elles font partie de l'historique codé. Une numérotation primitive récursive des coordonnées finies d'un état étant fixée, tout historique inverse \(x_\infty=(x_N)_{N\geq1}\in X_\infty^{\rm enr}\) détermine le code
\[
 \operatorname{code}(x_\infty)
 :=\{\langle N,j,c\rangle:(x_N)_j\text{ a pour code }c\}\subseteq\omega.
\]
Les équations de cohérence
\(\rho_{N+1,N}(x_{N+1})=x_N\) permettent de décoder de manière unique
l’histoire à partir de ce réel. Nous écrivons \(m_\infty\) pour le code de
l’histoire réalisée et \(m_\infty(n)\) pour sa fonction caractéristique sous la
numérotation fixée. Le code opératoire \(h\) et ce registre complet sont
réunis, au moyen d’une fonction primitive récursive d’appariement, dans le réel

\[
 u_{h,m}:=
 \{\langle0,\langle n,h(n)\rangle\rangle:n<\omega\}
 \cup
 \{\langle1,\langle n,m_\infty(n)\rangle\rangle:n<\omega\}.
\tag{33a}\label{hmtfg:base:eq:ch-codigo-conjunto}
\]

Les deux projections primitives récupèrent \(h\) et \(m_\infty\) à partir de \(u_{h,m}\).

La première extension est le prolongement métrique et opératoriel des états finis. Pour \(N\ge1\), soit \(\widehat{\mathcal X}^{\rm enr,raw}_{6N}\) l'ensemble des états enrichis bruts de longueur \(6N\), et définissons en termes ensemblistes le sous-système survivant par
\[
\begin{aligned}
 X_N
 &:=\widehat{\mathcal X}^{\rm enr,surv}_{6N}\\
 &:=\Bigl\{x\in\widehat{\mathcal X}^{\rm enr,raw}_{6N}:\\[-.2ex]
 &\qquad\forall M>N\ \exists z\in
   \widehat{\mathcal X}^{\rm enr,raw}_{6M}\;\text{ tel que}\\[-.2ex]
 &\qquad\operatorname{tr}_{6M,6N}(z)=x\ \wedge\
   z\text{ satisfait chaque transition TPK intermédiaire}\Bigr\}.
\end{aligned}
\]
et posons
\[
 \rho_{N+1,N}:=\operatorname{tr}_{6N+6,6N}\!\upharpoonright X_{N+1},
\]
et
\[
 \operatorname{Ext}_N(x):=
 \{y\in X_{N+1}:(x,y)\in
 \operatorname{Ext}_{6N+6,6N}\}.
\]
Le niveau racine est inclus sans exception implicite :
\[
 X_0:=\{\ast\},\qquad
 \rho_{1,0}:X_1\to X_0\text{ est constant},\qquad
 \operatorname{Ext}_0(\ast):=X_1.
\]
L’exposant \(\mathrm{surv}\) désigne
l’arbre élagué des états qui possèdent des descendants TPK compatibles à tout
horizon fini. La loi de prolongement TPK établie dans les sections
précédentes prouve que les germes
émis appartiennent à cet arbre ; la ramification finie et le lemme de Kőnig
transforment la compatibilité à tout horizon en une branche infinie. En
particulier,

\[
 \forall N\ge0\;\forall x\in X_N\quad
 \varnothing\ne\operatorname{Ext}_N(x)
 \subseteq\rho_{N+1,N}^{-1}(x).
\tag{34}\label{hmtfg:base:eq:ch-extension-no-vacia}
\]

Par le même lemme de l’arbre à ramification finie, tout germe survivant
possède un prolongement compatible

\[
 x_\infty=(x_N)_{N\ge1}\in
 X_\infty^{\rm enr}:=\varprojlim_N(X_N,\rho_{N+1,N}).
\tag{35}\label{hmtfg:base:eq:ch-limite-inverso}
\]

La seconde extension agit \textbf{ensuite} sur le code complet de cette
histoire réalisée.
Elle ne s’infère pas de la seule existence de l’opérateur fini
\(\operatorname{Ext}_N\) : elle formalise à l’échelle ensembliste le principe HMT selon
lequel seuls sont objets du domaine intégral ceux qui possèdent une dérivation
généalogique depuis l’état et ses règles. L’extension sémantique transfinie
est :

\[
 \begin{aligned}
  \mathfrak G_0(h,m_\infty)
  &:=\operatorname{TC}\bigl(\{\omega,u_{h,m},h,m_\infty\}\bigr),\\
  \mathfrak G_{\beta+1}(h,m_\infty)
  &:=\operatorname{Def}_{\Sigma_{\rm HMT}}
       \bigl(\mathfrak G_\beta(h,m_\infty)\bigr),\\
  \mathfrak G_\lambda(h,m_\infty)
  &:=\bigcup_{\beta<\lambda}\mathfrak G_\beta(h,m_\infty)
       \quad(\lambda\text{ limite}),\\
  \mathfrak G_{\rm HMT}(h,m_\infty)
  &:=\bigcup_{\beta\in\operatorname{Ord}}
       \mathfrak G_\beta(h,m_\infty).
 \end{aligned}
\tag{36}\label{hmtfg:base:eq:ch-clausura-transfinita}
\]

L'hypothèse du continu et une bijection cible ne sont pas introduites parmi les conditions qui définissent l'opérateur successeur ; l'opérateur est donné prospectivement par

\[
 \operatorname{Def}_{\Sigma_{\rm HMT}}(M)
 :=M\cup
 \Bigl\{
   \{x\in M:\langle M,\in,
       (\dot R_S\cap M^{k_S})_{S\in\Sigma_{\rm HMT}}\rangle
       \models\varphi(x,\bar p)\}:
   \ulcorner\varphi\urcorner\in\omega,
   \ \bar p\in M^{<\omega}
 \Bigr\}.
\tag{37}\label{hmtfg:base:eq:ch-operador-def}
\]

Autrement dit, elle exprime simultanément tous les sous-ensembles définissables par une formule finie du langage relationnel HMT et des paramètres déjà produits par la généalogie. Elle ne consulte ni valeur réelle cible, ni constante conventionnelle, ni cardinal cible, ni modèle extérieur. La formule~\eqref{hmtfg:base:eq:ch-operador-def} définit la clôture transfinie sur les germes, les opérateurs, l'état enrichi, le registre et la double projection construits précédemment.

\begin{lemma}[Codage conjoint du programme HMT et du registre réalisé]
\label{hmtfg:base:thm:ch-6}
La paire formée par
le code opératoire \(h\) et le registre \(m_\infty\) est codée par un unique
réel \(u_{h,m}\subseteq\omega\), et toute formule de \(\Sigma_{\rm HMT}\) se
traduit uniformément en une formule du langage \(\{\in,u_{h,m}\}\).

\end{lemma}
\begin{proof}
\(h\) et \(m_\infty\) sont des suites d'entiers naturels : \(h\) code la signature, les tables finies et les programmes des opérateurs ; \(m_\infty\) code un seul historique réalisé, et non l'espace complet des historiques. La formule~\eqref{hmtfg:base:eq:ch-codigo-conjunto} et ses deux décodeurs prouvent la première affirmation. Chaque symbole de~\eqref{hmtfg:base:eq:ch-signatura} désigne sa \textbf{relation de graphe codée dans \(h\)}, jamais le graphe externe de toutes ses évaluations possibles. APP, TRIT, \(U_t\), \(\Gamma_9\), les restrictions \(\rho\) et les extensions finies ont des graphes primitifs récursifs sur des codes entiers et, par conséquent, des formules \(\Delta_0\) absolues entre niveaux transitifs. Les deux représentations s'expriment comme relations de graphe au moyen des formules

\[
 \begin{aligned}
 \dot R_{\rm ar}(x,y)
 &\;\Longleftrightarrow\;
 \forall k\in\omega\;\exists N\;\forall n\ge N\quad
 |\mu_n(x)-y|<2^{-k},\\
 \dot R_{\rm exc}(x,c,z)
 &\;\Longleftrightarrow\;
 \forall N\in\omega\quad z\!\upharpoonright N
   =E_N(x\!\upharpoonright N,c\!\upharpoonright N),
 \end{aligned}
\tag{37c$_0$}\label{hmtfg:base:eq:ch-grafos-publicacion}
\]

où \(\mu_n\) et \(E_N\) sont les fonctions rationnelles/finies dont les codes
figurent dans \(h\). On remplace par récurrence chaque formule atomique contenant
un symbole HMT par sa formule de graphe de paramètre \(u_{h,m}\). Les
connecteurs et les quantificateurs sont conservés. Cette récurrence sur la complexité
syntaxique produit une traduction \(\varphi\mapsto\varphi^*\) telle que, pour
tout niveau transitif \(M\) de~\eqref{hmtfg:base:eq:ch-clausura-transfinita} contenant les paramètres,

\[
 \langle M,\in,(\dot R_S\cap M^{k_S})_S\rangle
   \models\varphi(\bar a)
 \quad\Longleftrightarrow\quad
 M\models\varphi^*(\bar a;u_{h,m}).
\tag{37c}\label{hmtfg:base:eq:ch-eliminacion-uniforme}
\]

Ainsi, \eqref{hmtfg:base:eq:ch-operador-def} ne peut consulter ni oracle, ni table décimale cible, ni relation non codée introduisant une infinité non dénombrable d'objets en une seule étape.
\end{proof}

La sémantique HMT est maintenant définie \textbf{indépendamment de son identification à \(\mathfrak G\)}. Soit \(\mathcal T_0^{\rm HMT}\) l'ensemble des noms canoniques des éléments de \(\mathfrak G_0(h,m_\infty)=
\operatorname{TC}(\{\omega,u_{h,m},h,m_\infty\})\). En particulier, le niveau initial contient le registre réalisé \(m_\infty\), et non toutes les histoires de \(X_\infty^{\rm enr}\) comme paramètres primitifs. Les termes admissibles sont les arbres bien fondés obtenus par

\[
 \begin{aligned}
 \mathcal T^{\rm HMT}_{\beta+1}
 &:={\mathcal T}^{\rm HMT}_\beta\cup
 \{\langle\beta,\ulcorner\varphi\urcorner,
       t_0,\ldots,t_{k-1}\rangle:
       t_i\in\mathcal T^{\rm HMT}_\beta\},\\
 \mathcal T^{\rm HMT}_\lambda&:=
       \bigcup_{\beta<\lambda}\mathcal T^{\rm HMT}_\beta,
 \qquad
 \mathcal T^{\rm HMT}:=
       \bigcup_{\beta\in\operatorname{Ord}}\mathcal T^{\rm HMT}_\beta.
 \end{aligned}
\tag{37a}\label{hmtfg:base:eq:ch-terminos}
\]

La valeur d’un terme successeur est le sous-ensemble du niveau précédent défini
par \(\varphi\) dans la structure relationnelle induite, avec les valeurs de
\(t_0,\ldots,t_{k-1}\) comme paramètres. On définit, avant toute
comparaison ensembliste,

\[
 \operatorname{Ob}_{\rm sem}^{\rm HMT}(h,m_\infty)
 :=\{\llbracket t\rrbracket:t\in\mathcal T^{\rm HMT}\}.
\tag{37b}\label{hmtfg:base:eq:ch-objetos-semanticos}
\]

C'est la règle d'\textbf{exhaustivité générative} adoptée par HMT : tout objet du domaine intégral doit posséder l'un de ces arbres de dérivation, et tout arbre admissible exprime un objet. Elle n'est pas une conséquence de la ramification finie. L'hypothèse du continu n'est pas introduite dans cette règle ; elle est obtenue par la démonstration développée dans l'article. Ni une énumération des réels ni une bijection cible ne sont introduites comme données.

\begin{theorem}[Équivalence entre la sémantique des termes et la hiérarchie généalogique]
\label{hmtfg:base:thm:ch-6a}
L’évaluation des
termes précédents vérifie

\[
 \operatorname{Ob}_{\rm sem}^{\rm HMT}(h,m_\infty)
 =\mathfrak G_{\rm HMT}(h,m_\infty).
\tag{37d}\label{hmtfg:base:eq:ch-equivalencia-semantica}
\]

\end{theorem}
\begin{proof}
Par récurrence transfinie. Au niveau initial, les deux classes coïncident. Si toute valeur d'un terme de \(\mathcal T^{\rm HMT}_\beta\) appartient à \(\mathfrak G_\beta\), la sémantique du terme successeur est l'un des sous-ensembles de~\eqref{hmtfg:base:eq:ch-operador-def}. Réciproquement, chaque sous-ensemble de~\eqref{hmtfg:base:eq:ch-operador-def} est donné par un code de formule et un tuple fini de paramètres ; d'après l'hypothèse de récurrence, ces paramètres possèdent des termes et~\eqref{hmtfg:base:eq:ch-terminos} forme le terme qui le désigne. À une limite, les deux constructions prennent la même réunion.
\end{proof}

\Needspace{9\baselineskip}
L’encodeur d’états est compatible avec les troncatures et envoie chaque
prolongement TPK sur un terme successeur. Soit
\(\mathcal T_N^{\rm st}\subset\mathcal T^{\rm HMT}\) la sous-classe de termes
qui code des états complets de profondeur \(N\), et soit
\(\tau_{N+1,N}:\mathcal T_{N+1}^{\rm st}\to\mathcal T_N^{\rm st}\) l’effacement
du dernier bloc avec la troncature correspondante du registre généalogique,
de la retenue et de la mémoire. Si \(J_N:X_N\to\mathcal T_N^{\rm st}\) est l’encodeur
complet, la compatibilité de \(\operatorname{Ext}_N\) prouve l’identité
correctement typée

\[
 \forall x\in X_N\ \forall y\in\operatorname{Ext}_N(x)\qquad
 \tau_{N+1,N}\bigl(J_{N+1}(y)\bigr)=J_N(x).
\tag{37e$_0$}\label{hmtfg:base:eq:ch-naturalidad-codificador}
\]

Définissons au sein de la sémantique
\[
 \operatorname{Succ}^{\rm TPK}_{N+1}(J_N(x)):=
 \{t\in\mathcal T_{N+1}^{\rm st}:\exists y\in X_{N+1}\,
   (y\in\operatorname{Ext}_N(x)\ \wedge\ t=J_{N+1}(y))\}.
\]
Choisissons \(\beta\) tel que \(\mathfrak G_\beta\) contienne \(x\), \(X_N\),
\(X_{N+1}\), les graphes finis de
\(\operatorname{Ext}_N,J_N,J_{N+1}\) et leurs codes dans \(h\). Alors la
formule primitive récursive du graphe de \(\operatorname{Ext}_N\) donne

\[
 J_{N+1}[\operatorname{Ext}_N(x)]
 =\operatorname{Succ}^{\rm TPK}_{N+1}(J_N(x)),
 \qquad
 J_{N+1}[\operatorname{Ext}_N(x)]
 \in\mathfrak G_{\beta+1},
\tag{37e}\label{hmtfg:base:eq:ch-sucesores-semanticos}
\]

et les deux membres sont formés par la formule de graphe de la même extension
TPK. Les équations~\eqref{hmtfg:base:eq:ch-naturalidad-codificador}--\eqref{hmtfg:base:eq:ch-sucesores-semanticos} sont le lien sémantique explicite : la
sous-classe des termes successeurs est exactement l’image des prolongements
admissibles et leur troncature commute. On n’identifie pas la fibre complète de
\(\rho\) à \(\operatorname{Ext}_N\), et l’on n’identifie pas
\(\operatorname{Def}\) à TPK par nomenclature.

\begin{lemma}[Transitivité, axiomes et bon ordre du modèle généalogique]
\label{hmtfg:base:thm:ch-6b}
\(\mathfrak G_{\rm HMT}(h,m_\infty)\) est une classe transitive, contient tous les ordinaux, satisfait la théorie des ensembles de Zermelo--Fraenkel (ZF) et possède un bon ordre global définissable ; elle satisfait donc le choix et tous les axiomes utilisés dans la comparaison cardinale ultérieure.

\end{lemma}
\begin{proof}
Abrégeons \(B=\mathfrak G_0\) et \(u=u_{h,m}\). La récurrence sur~\eqref{hmtfg:base:eq:ch-clausura-transfinita} prouve simultanément que chaque niveau est transitif et que \(\mathfrak G_\alpha\in\mathfrak G_{\alpha+1}\) ; si tous les ordinaux inférieurs à \(\alpha\) sont présents, leur ensemble est définissable au niveau correspondant et exprime \(\alpha\). L'union contient donc tous les ordinaux. L'extensionnalité et la fondation sont héritées de l'univers ambiant ; l'ensemble vide et l'infini appartiennent à \(B\). Un niveau contenant les paramètres étant choisi, les formules définissant \(\{a,b\}\), \(\bigcup a\) et les sous-ensembles par séparation les expriment à un niveau successeur.

Pour la collection--remplacement, soit \(F\) une relation fonctionnelle définissable dans \(\mathfrak G\) sur l'ensemble \(a\). Le schéma de réflexion relatif à la hiérarchie définissable \((\mathfrak G_\xi,\in,u_{h,m})_{\xi\in\mathrm{Ord}}\), appliqué dans la métathéorie ambiante à la liste finie de formules qui définit \(F\), produit un niveau \(\mathfrak G_\theta\) qui contient \(a\), les paramètres, et qui est correct pour ces formules. Les rangs des témoins de \(F[a]\) sont alors bornés par \(\theta\) ; la séparation sur \(\mathfrak G_\theta\) forme l'image dans \(\mathfrak G_{\theta+1}\). Pour l'ensemble des parties, on ne présuppose pas l'ensemble interne que l'on veut construire. Une fois \(a\in\mathfrak G_\gamma\) fixé, on considère dans la métathéorie ambiante l'ensemble déjà existant \(\mathcal P^V(a)\). Pour chaque \(b\in\mathcal P^V(a)\cap\mathfrak G\), soit \(r(b)\) sa première étape généalogique. Comme \(\mathfrak G\) est une classe définissable, la séparation dans la métathéorie ambiante forme d'abord l'ensemble \(\mathcal P^V(a)\cap\mathfrak G\). Le remplacement \textbf{dans la métathéorie ambiante} appliqué à cet ensemble borne les ordinaux \(r(b)\) par un certain \(\theta\). Par conséquent,
\[
 \mathcal P^{\mathfrak G}(a)
 =\{b\in\mathfrak G_\theta:b\subseteq a\},
\]
et le membre droit s'exprime par séparation dans \(\mathfrak G_{\theta+1}\). Il n'introduit pas dans \(\mathfrak G\) les sous-ensembles extérieurs à la sémantique~\eqref{hmtfg:base:eq:ch-objetos-semanticos}. C'est l'argument non circulaire de majoration par les rangs.

Enfin, on attribue à chaque objet son premier niveau de naissance et son
plus petit \textbf{nom de dérivation} : un arbre de dérivation bien fondé et
à ramifications finies (et, par~\eqref{hmtfg:base:eq:ch-terminos}, fini pour chaque nom individuel), étiqueté par un
ordinal de niveau, un code de formule et le tuple des noms de ses paramètres.
Par récurrence transfinie, on ordonne ces noms : d’abord l’étiquette ordinale,
puis le code de formule, puis, lexicographiquement, les noms des
paramètres déjà ordonnés. Le plus petit nom de chaque valeur définit une relation de
classe \(<_{\mathfrak G}\) qui compare tous les objets. C’est un bon ordre
global définissable qui, en prenant le plus petit élément de chaque ensemble non vide,
prouve le Choix.
\end{proof}

% Propietario reproducido por unidad demostrativa; véase PROPIETARIOS_CONTINUO_UNIVERSO.json.
% Origen: XI ES CUMPLIMIENTO_EDITORIAL_20260928, sections/ch_reflexion_testigos.tex:1--89.
% Cuerpos y pruebas conservados; sólo namespace, rutas y remisiones contextuales declaradas.
\paragraph{Réflexion par clôture des témoins et construction de l'image.}
L’étape de Collection--Remplacement du lemme~\ref{hmtfg:base:thm:ch-6b} admet la
démonstration explicite suivante. On maintient la hiérarchie déjà engendrée à partir du
code opératoire et du registre réalisé ; nous abrégeons
\(\mathfrak G=\mathfrak G_{\rm HMT}(h,m_\infty)\).
L’argumentation utilise la Séparation et le Remplacement dans la métathéorie ambiante,
sans présupposer ces schémas dans \(\mathfrak G\).

Soit \(\Phi\) une collection finie de formules. On traduit d'abord leurs quantificateurs universels au moyen de la négation et de la quantification existentielle, puis on ferme la collection résultante par sous-formules. Pour tout niveau initial, il existe un ordinal limite ultérieur \(\theta\) tel que
\[
 \mathfrak G_\theta\models\varphi(\bar a)
 \quad\Longleftrightarrow\quad
 \mathfrak G\models\varphi(\bar a)
 \qquad
 (\varphi\in\Phi,\ \bar a\in\mathfrak G_\theta^{<\omega}).
\]
La notation de satisfaction de la classe est interprétée, pour chaque formule fixée, en relativisant ses quantificateurs à la classe définissable \(\mathfrak G\) ; aucun prédicat uniforme de vérité n'est introduit.

\begin{proof}
Fixons \(\alpha\). Pour chaque sous-formule existentielle
\(\exists y\,\psi(y,\bar x)\) de \(\Phi\), les tuples
\(\bar a\in\mathfrak G_\alpha^{<\omega}\) de l’arité correspondante
qui possèdent un témoin dans \(\mathfrak G\) forment un ensemble par Séparation
métathéorique. À chacun est attribué le plus petit ordinal \(\beta\) pour lequel
un \(y\in\mathfrak G_\beta\) satisfait
\(\mathfrak G\models\psi(y,\bar a)\). L’ordinal existe car
\(\mathfrak G\) est l’union de ses niveaux. Le Remplacement dans la métathéorie
borne ces premières étapes des témoins. La finitude de \(\Phi\) permet
de choisir de manière définissable une seule borne \(f(\alpha)>\alpha\), valable pour
toutes ses sous-formules existentielles. Par conséquent, chaque tuple considéré possède
\emph{au moins un témoin} dans \(\mathfrak G_{f(\alpha)}\).

Nous prenons \(\alpha_0\) suffisamment grand pour contenir les paramètres
initiaux et posons
\[
 \alpha_{n+1}=f(\alpha_n),\qquad
 \theta=\sup_{n<\omega}\alpha_n.
\]
Par continuité de la hiérarchie, tout tuple fini de
\(\mathfrak G_\theta\) appartient à un certain \(\mathfrak G_{\alpha_n}\).
Si une formule existentielle de \(\Phi\) avec ce tuple est vraie dans
\(\mathfrak G\), elle possède un témoin dans
\(\mathfrak G_{\alpha_{n+1}}\subseteq\mathfrak G_\theta\).
La récurrence sur les formules prouve maintenant l’équivalence : les formules atomiques
sont conservées dans la structure induite ; les connecteurs préservent
l’équivalence ; l’étape existentielle utilise ce témoin dans un sens et le
même témoin, vu dans \(\mathfrak G\), dans l’autre. La traduction préalable
des universelles inclut aussi les formules existentielles qui expriment
leurs contre-exemples possibles. La réflexion finie requise est prouvée.
\end{proof}

Soit maintenant \(F(x,y,\bar p)\) fonctionnelle sur \(a\) dans \(\mathfrak G\). Nous appliquons la construction aux formules qui expriment \(F\), l'existence de ses témoins et son image, avec \(a,\bar p\in\mathfrak G_\theta\). La transitivité donne \(a\subseteq\mathfrak G_\theta\). Chaque \(x\in a\) possède sa valeur à l'intérieur de ce niveau, et la réflexion identifie l'image exacte à
\[
 b=\{y\in\mathfrak G_\theta:
       \mathfrak G_\theta\models
       \exists x\in a\ F(x,y,\bar p)\}.
\]
Cet ensemble est définissable sur \(\mathfrak G_\theta\), avec des paramètres du même niveau. Par l'opération successeur de la hiérarchie, \(b\in\operatorname{Def}(\mathfrak G_\theta)
\subseteq\mathfrak G_{\theta+1}\). On obtient ainsi le remplacement dans \(\mathfrak G\). Si l'on conserve l'existence de témoins et que l'on omet l'unicité, la même construction fournit un ensemble de témoins pour la collection.

\paragraph{Arité des noms et sélection interne.}
Dans le bon ordre du lemme~\ref{hmtfg:base:thm:ch-6b}, l’énumération de la base
attribue à chaque élément son plus petit code naturel. À chaque successeur, on ordonne
d’abord les codes de formule ; une fois l’un fixé, l’arité de son
tuple de paramètres est fixée. L’ordre lexicographique s’applique alors à un
\emph{produit fini d’arité fixe} des bons ordres précédents.
Le plus petit nom représente chaque nouvelle valeur et les niveaux limites conservent
l’ordre de première apparition. L’ordre des tuples est ainsi spécifié,
sans mélanger des arités différentes dans un unique ordre
lexicographique. La sélection du plus petit élément de chaque ensemble non vide,
avec le Remplacement déjà obtenu, forme intérieurement les fonctions de
choix utilisées ensuite. Cette construction de noms précède les
injections cardinales et ne reçoit aucune bijection du continu comme
donnée.

